\begin{frame}{QuIM-RAG: Advancing Retrieval-Augmented Generation With Inverted Question Matching}
\scriptsize
\vspace{-0.8em}

\begin{table}[h!]
\centering
\renewcommand{\arraystretch}{1.5} % increases row height slightly more
\setlength{\tabcolsep}{2.5pt}     % reduces horizontal padding a bit more

\begin{tabular}{|p{2.6cm}|p{4.3cm}|p{4.3cm}|}
\hline
\textbf{Citation} & \textbf{Summary} & \textbf{Advantages \& Disadvantages} \\
\hline
\tiny\textit{[5] B. Saha, U. Saha, and M. Z. Malik, IEEE Access, vol. 12, pp. 185401–185410, 2024.  
DOI: 10.1109/ACCESS.2024.3513155.}
&
Introduces \textbf{QuIM-RAG}, a Retrieval-Augmented Generation model using \textbf{question-to-question inverted index matching} with quantized embeddings to improve QA relevance. Tested on NDSU data using Meta-LLaMA3-8B-instruct.It generates candidate questions from document chunks to retrieve more semantically relevant answers and outperform traditional RAG approaches.
&
\textbf{Advantages:}
\begin{itemize}\itemsep0.25em
\item More semantically relevant retrieval.
\item High faithfulness score (1.00).
\item Strong BERT/RAGAS performance.
\end{itemize}

\textbf{Disadvantages:}
\begin{itemize}\itemsep0.25em
\item Evaluation limited to a single domain.
\item High preprocessing and storage cost.
\item Dependent on quality of generated questions.
\end{itemize} \\
\hline
\end{tabular}
\end{table}
\end{frame}
